%=============================================================================
% preamble.tex
% Arcaea共通テスト（ACT）テンプレート共通プリアンブル
%
% 使用方法：
%   
% 2. 基本グラフィック・色・レイアウト
\usepackage{fancybox}
\usepackage{graphicx,color}
\usepackage[margin=20mm]{geometry}
\usepackage{multicol,multirow,varwidth,float,caption}

% 3. emathコアパッケージ
\usepackage[deepenum,notMy]{emath}

% emath拡張パッケージ群
\usepackage{hako}        % 穴埋め（箱）
\usepackage{EMfbox}      % 囲み枠
\usepackage{abcgreek}    % ギリシャ文字
\usepackage{emathEy}     % 記号・特殊文字
\usepackage{emathG}      % 幾何（図形）
\usepackage{emathT}      % 表作成
\usepackage{emathPs}     % PostScript連携
\usepackage{emathR}      % 実数計算
\usepackage{emathAe}     % 解答作成
\usepackage{emathThm}    % 定理環境
\usepackage{EMpict2e}    % 描画拡張
\usepackage[deluxe]{emathOtf} % OTFフォント（太字などの対応）


% 5. 各種定義と書式設定
% 長丸囲み文字の定義
\def\nagamaru#1{\mbox{\pIIebNagamaru{#1}}}

% 定理環境（「仮定」など）の共通テスト風デザイン
\theorembodyfont{\normalfont}
\theoremstyle{boxed}
\newtheorem<frame=rectbox,frameoption={oval=8pt,preitem=~,postitem=~}>{katei}{仮定}

% キャプションの日本語化（図1: → 図 1）
\captionsetup{figurename=図,labelsep=quad}

% カウンターのリセット設定（大問ごとに数式や図の番号をリセット）
\resetcounter{equation}[enumi]
\resetcounter{figure}[enumi]
\resetcounter{table}[enumi]
\EMfigtrue
\apnlist{\setlength{\listparindent}{1zw}}

% 大問ラベルの書式定義（専用の箇条書き環境として定義）
\newenvironment{QuestionEnum}{%
    \begin{enumerate}<apnlist={\leftmargin1zw\itemindent4zw}>'{\bfseries\Large 第\arabic{enumi}問}'%
}{%
    \end{enumerate}%
}


% --- 単線枠 ---
\newcommand{\fgb}[1]{%
    {% 
    \setlength{\fboxrule}{1.6pt}%
    \framebox[1.7cm]{#1}%
    }%
}
\newcommand{\fb}[1]{%
    {% 
    \setlength{\fboxrule}{0.4pt}%
    \framebox[1.7cm]{#1}%
    }%
}

% --- 二重線枠 ---
\newcommand{\dgb}[1]{%
    {%
    \edef\originalsep{\the\fboxsep}% 現在の標準の余白サイズを記憶
    \setlength{\fboxrule}{1.6pt}% ① 外側の線を太くする(1.6pt)
    \setlength{\fboxsep}{1.6pt}%  ② 二重線の「線と線のすき間」の幅
    \fbox{%
        \setlength{\fboxrule}{0.4pt}% ③ 内側の線は普通の太さ(0.8pt)
        \fbox{\makebox[1.4cm]{#1}}%
    }%
    }%
}
\newcommand{\db}[1]{%
    {%
    \edef\originalsep{\the\fboxsep}% 現在の標準の余白サイズを記憶
    \setlength{\fboxrule}{0.4pt}% ① 外側の線を太くする(1.6pt)
    \setlength{\fboxsep}{1.6pt}%  ② 二重線の「線と線のすき間」の幅
    \fbox{%
        \setlength{\fboxrule}{0.4pt}% ③ 内側の線は普通の太さ(0.8pt)
        \fbox{\makebox[1.4cm]{#1}}%
    }%
    }%
}

\newcommand{\subject}[3]{%
\begin{center}
    {\LARGE \textbf{#1}}\\[0.5em]
    {\Huge （}%
    \raisebox{0.4em}{%
        {\normalsize \textbf{解答番号}} \fb{#2} ～ \fb{#3}%
    }%
    {\Huge ）}\\
\end{center}
}  を ACTXXXX_main.tex のプリアンブルに記述する
%
% 依存パッケージ：
%   emath（TeX Live 標準外・別途インストール必要）
%   その他はTeX Live標準同梱
%=============================================================================


\usepackage{fancybox}       % \doublebox などの枠線コマンド（emath と別系統）
\usepackage{graphicx,color} % 画像挿入・色指定（\includegraphics, \textcolor 等）
\usepackage[margin=20mm]{geometry}
                            % ページ余白設定（上下左右すべて20mm）
\usepackage{multicol,multirow,varwidth,float,caption}
                            % multicol  : 2段組（選択肢の横並びなど）
                            % multirow  : 表のセル結合（注意事項の配点表等）
                            % varwidth  : 中身に合わせた幅のminipage
                            % float     % 図表の配置制御（[H]オプション等）
                            % caption   : キャプション書式変更（figurename=図 等）

\usepackage[deepenum,notMy]{emath}
                            % emath 本体（必須）
                            % deepenum : enumerate を深いレベルまで使用可能にする
                            % notMy    : emath 独自の一部定義を抑制（他パッケージとの衝突回避）
\usepackage{hako}           % 解答欄（穴埋め箱）: \fgb{}, \fb{}, \dgb{}, \db{} 等
\usepackage{EMfbox}         % 枠付きボックス: \rectbox, \EMbox 等
\usepackage{abcgreek}       % ギリシャ文字アルファベット対応（\alpha 等の表示改善）
\usepackage{emathEy}        % 選択肢の横並び列挙環境（\yokoenum 等）
\usepackage{emathG}         % 幾何作図サポート（座標平面の描画等）
\usepackage{emathT}         % 表作成拡張（増減表など）
\usepackage{emathPs}        % PostScript 連携描画（\zahyou 等の座標系グラフ）
\usepackage{emathR}         % 外部ソースファイルの読み込み（\ReadTeXFile の依存先）
\usepackage{emathAe}        % 解答ファイル自動生成機能
                            % \openKaiFile / \closeKaiFile で解答ファイルを開閉
                            % 問題ファイル内の \解答{} で内容を記録
                            % \ReadTeXFile もこのパッケージが提供する
\usepackage{emathThm}       % 定理・仮定環境（\theoremstyle{boxed} 等）
\usepackage{EMpict2e}       % pict2e 拡張（emath の描画機能を pict2e で補完）
\usepackage[deluxe]{emathOtf}
                            % OTFフォント対応（太字ゴシック等）
                            % \textgt{} による明朝/ゴシック切り替えが可能になる

% 長丸囲み文字
% \nagamaru{テキスト} で長丸で囲んだ文字を出力する
\def\nagamaru#1{\mbox{\pIIebNagamaru{#1}}}

% 定理環境（「仮定」ブロック）の共通テスト風デザイン
% 角丸の枠付き定理環境「仮定」を定義する
% 問題文中で条件・前提をまとめるブロックとして使用
\theorembodyfont{\normalfont}
\theoremstyle{boxed}
\newtheorem<frame=rectbox,frameoption={oval=8pt,preitem=~,postitem=~}>{katei}{仮定}

% キャプションの日本語化
% 図のキャプション「Figure 1:」を「図 1　」形式に変更
\captionsetup{figurename=図,labelsep=quad}

% カウンターのリセット設定
% 大問（enumerate第1レベル）が変わるたびに数式番号・図番号・表番号をリセット
\resetcounter{equation}[enumi]
\resetcounter{figure}[enumi]
\resetcounter{table}[enumi]
\EMfigtrue
\apnlist{\setlength{\listparindent}{1zw}}



% 大問番号を「第1問」「第2問」... と太字大文字で表示する enumerate 環境
\newenvironment{QuestionEnum}{%
    \begin{enumerate}<apnlist={\leftmargin1zw\itemindent4zw}>'{\bfseries\Large 第\arabic{enumi}問}'%
}{%
    \end{enumerate}%
}

% 数値・記号穴埋め欄 : \fgb{} / \fb{}
% 選択式解答欄       : \dgb{} / \db{}

% 数値・記号穴埋め欄（太枠）: 初出用
\newcommand{\fgb}[1]{%
    {%
    \setlength{\fboxrule}{1.6pt}%   外枠を太く（1.6pt）
    \framebox[1.7cm]{#1}%
    }%
}

% 数値・記号穴埋め欄（細枠）: 2回目以降用
\newcommand{\fb}[1]{%
    {%
    \setlength{\fboxrule}{0.4pt}%   外枠を通常の細さ（0.4pt）
    \framebox[1.7cm]{#1}%
    }%
}

% 選択式解答欄（太枠二重）: 初出用
\newcommand{\dgb}[1]{%
    {%
    \edef\originalsep{\the\fboxsep}%
    \setlength{\fboxrule}{1.6pt}%   外側の線を太く（1.6pt）
    \setlength{\fboxsep}{1.6pt}%    二重線の線間隔
    \fbox{%
        \setlength{\fboxrule}{0.4pt}% 内側の線は通常の細さ（0.4pt）
        \fbox{\makebox[1.4cm]{#1}}%
    }%
    }%
}

% 選択式解答欄（細枠二重）: 2回目以降用
\newcommand{\db}[1]{%
    {%
    \edef\originalsep{\the\fboxsep}%
    \setlength{\fboxrule}{0.4pt}%   外側の線を細く（0.4pt）
    \setlength{\fboxsep}{1.6pt}%    二重線の線間隔
    \fbox{%
        \setlength{\fboxrule}{0.4pt}%
        \fbox{\makebox[1.4cm]{#1}}%
    }%
    }%
}

% \subject{科目名}{開始解答番号}{終了解答番号}
\newcommand{\subject}[3]{%
\begin{center}
    {\LARGE \textbf{#1}}\\[0.5em]
    {\Huge （}%
    \raisebox{0.4em}{%
        {\normalsize \textbf{解答番号}} \fb{#2} ～ \fb{#3}%
    }%
    {\Huge ）}\\
\end{center}
}